\chapter{Conclusions and Future Work} 
\label{ch:conclusion} 
\todo{inline}{WIP}
\section{Summary of Thesis} \label{ch7:summary}
This thesis focuses on the development of domain-specific datasets, models, and evaluation frameworks tailored to the challenges of online discourse. We examined the limitations of applying general-purpose language models to social media, highlighting the role of informality, multilinguality, and rapid linguistic change. Our work improves on limitations on  core \ac{NLP} tasks specifically topic classification and hate speech detection as well as providing a robust benchamrk for evaualting models on social media content. Furthermore, we showcase the value of these resources by utilising them to perform in depth analysis on topics such as misinformation and polarisation by establishing clear and robust methodologies. The research was structured to first create new datasets and benchmarks, then apply them to downstream analyses, producing insights into user behaviour, discourse patterns, and the spread of harmful content online.

\section{Contributions}
\label{ch7:contributions}

The thesis makes several contributions to \ac{NLP} for social media. We introduced new datasets for topic classification, including a multilingual extension, filling the gap in high-quality, in-domain resources for this task. We developed and evaluated robust hate speech detection models through large-scale cross-dataset experiments, showing how combining resources can improve generalisability. We proposed a unified benchmark, {\DATASET}, for evaluating language models across a diverse set of social media–specific tasks, enabling systematic comparison of general-purpose and domain-adapted models. 

Additionaly, we applied these resources and models to investigate misinformation spreaders and polarised discourse, offering new perspectives on the linguistic and behavioural features that shape online debates. While these steps represent substantial progress, there is room to build on this work by extending ......



\section{Future Work} \label{ch7:future-work}
The datasets introduced here could be expanded to include additional languages, domains. The unified benchmark can be expanded with new tasks and evaluation metrics.....